% Original flowchart for the owner's paper-reproduce workflow.
% Content: assets/decision-stages.mmd and the approved two-stage decisions.
% Renderer: LuaLaTeX/TikZ. Profile: schematic-designer Minimalist.
% Design lesson: aligned groups and reserved connector space, from the local
% lab-curriculum-flow tutorial. No example or upstream artwork/code copied.
% Palette: teal for objective; warm gold for depth. All alternatives within
% a stage share one color; stage labels retain the distinction without color.
% Geometry indicates choices and ordering only; there are no measured quantities.
% Generate minimalist-profile.sty from schematic-designer's shared profile.
\documentclass[tikz,border=2mm]{standalone}
\usepackage{minimalist-profile}
\usetikzlibrary{calc}
\colorlet{objectiveStroke}{msColor5!75!msInk}
\colorlet{objectiveText}{msColor6}
\colorlet{objectiveCard}{msColor5!8!msPaper}
\colorlet{objectiveGroup}{msColor5!3!msPaper}
\colorlet{objectiveBorder}{msColor5!35!msPaper}
\colorlet{depthStroke}{msColor3!60!msInk}
\colorlet{depthText}{msColor3!40!msInk}
\colorlet{depthCard}{msColor3!14!msPaper}
\colorlet{depthGroup}{msColor3!4!msPaper}
\colorlet{depthBorder}{msColor3!50!msPaper}
\begin{document}
\pagecolor{msPaper}
\hyphenpenalty=10000
\exhyphenpenalty=10000
\begin{tikzpicture}[ms base,x=1mm,y=1mm,
    stage/.style={draw=msMuted!60,fill=msPaper,line width=\msLineContext,rounded corners=1.5mm},
    card/.style={draw=msInk,fill=msPaper,line width=\msLineStructure,rounded corners=1.4mm,minimum width=52mm,minimum height=24mm,inner sep=0mm,outer sep=0mm},
    heading/.style={ms title,anchor=west,inner sep=0mm},
    card title/.style={ms title,anchor=north,inner sep=0mm,text width=46mm,align=center},
    description/.style={ms body,anchor=north,inner sep=0mm,text width=46mm,align=center}]

% Each stage is one group. No connector joins individual alternatives.
\draw[stage,draw=objectiveBorder,fill=objectiveGroup] (0,0) rectangle (176,-42);
\node[heading,text=objectiveText] at (5,-6) {Stage 1: What is my primary objective?};
\node[card,draw=objectiveStroke,fill=objectiveCard] (reproduction) at (30,-25) {};
\node[card,draw=objectiveStroke,fill=objectiveCard] (assessment) at (88,-25) {};
\node[card,draw=objectiveStroke,fill=objectiveCard] (application) at (146,-25) {};
\node[card title,text=objectiveText] at (30,-16) {Reproduction};
\node[card title,text=objectiveText] at (88,-16) {Assessment};
\node[card title,text=objectiveText] at (146,-16) {Application};
\node[description] at (30,-22) {Reproduce the paper's model, selected features, or reported results.};
\node[description] at (88,-22) {Verify or validate a model, test a mechanism or hypothesis, or assess agreement with evidence.};
\node[description] at (146,-22) {Use, adapt, or extend the paper's model, mechanism, or method for my own scientific question.};

\draw[ms arrow] (88,-42) -- (88,-56);

\draw[stage,draw=depthBorder,fill=depthGroup] (0,-56) rectangle (176,-98);
\node[heading,text=depthText] at (5,-62) {Stage 2: What depth of work do I need?};
\node[card,draw=depthStroke,fill=depthCard] (model) at (30,-81) {};
\node[card,draw=depthStroke,fill=depthCard] (claim) at (88,-81) {};
\node[card,draw=depthStroke,fill=depthCard] (result) at (146,-81) {};
\node[card title,text=depthText] at (30,-72) {Model Level};
\node[card title,text=depthText] at (88,-72) {Claim Level};
\node[card title,text=depthText] at (146,-72) {Result Level};
\node[description] at (30,-78) {Implement, adapt, or check the required computational model or method.};
\node[description] at (88,-78) {Test a selected scientific claim, proposed mechanism, or hypothesis.};
\node[description] at (146,-78) {Reproduce, compare, or generate specific curves, tables, or numerical results.};

\node[ms small,anchor=north,inner sep=0mm,text=msMuted] at (88,-103) {Choose one objective and one depth; any combination is allowed.};
\end{tikzpicture}
\end{document}
